\section{Methodology}
\label{sec:method}

This section decomposes the online control problem of
Section~\ref{sec:intro} into three sub-problems and solves each under
causal information alone. Section~\ref{sec:problem} formalizes the
two-tier pool and the cost asymmetry that makes victim choice the
dominant lever. Section~\ref{sec:esp} introduces \spear, which
approximates the B\'el\'ady eviction oracle from two observable
signals, namely a gap recurrence estimate and a turn-indexed
completion hazard. Section~\ref{sec:okt} introduces \tide, which converts
measured idle windows into tier-placement migrations ranked by the
same signal. Section~\ref{sec:loop} unifies both mechanisms in a
single event-driven control loop as shown in Algorithm~\ref{alg:unison}.
Fig.~\ref{fig:method} gives the corresponding pipeline.

\begin{figure*}[t]
\centering
\includegraphics[width=\textwidth]{fig02_methodology.pdf}
\caption{\unison methodology. A shared event log bridges workload
characterization, survival-based scoring, idle-window migration, and
multi-metric evaluation.}
\label{fig:method}
\end{figure*}

\subsection{Workload Model and Problem Formulation}
\label{sec:problem}

An agent execution is a session $s$ of $T_s$ turns. Turn $k$ arrives at
wall time $a_{s,k}$ after a tool gap
\begin{equation}
  g_{s,k}=a_{s,k}-c_{s,k-1},
  \label{eq:gap}
\end{equation}
where $c_{s,k-1}$ is the completion time of the previous turn.
The resident KV length $\ell_{s,k}$ is prefix-heavy because later turns
inherit most of their context from earlier ones, so a hit avoids
re-prefilling the full prefix while a miss pays the entire sequence
length. The reuse ratio
\begin{equation}
  \rho_{s,k}
  =\frac{\ell_{s,k-1}}{\ell_{s,k}},
  \qquad
  \Delta\ell_{s,k}
  =\ell_{s,k}-\ell_{s,k-1}
  \label{eq:reuse}
\end{equation}
quantifies this asymmetry. A miss costs
$\ell_{s,k}\,t_{\mathrm{p}}$ whereas a hit costs only
$\Delta\ell_{s,k}\,t_{\mathrm{p}}$, making victim choice the
dominant lever on serving latency.

The live set $\mathcal{L}(t)$ occupies a two-tier budget of SRAM
capacity $C_{\mathrm{S}}$ and HBM capacity $C_{\mathrm{H}}$ together
with a hard session-slot limit $N_{\mathrm{sess}}$,
\begin{equation}
  \sum_{s\in\mathcal{L}_{\mathrm{S}}}\ell_s\le C_{\mathrm{S}},
  \qquad
  \sum_{s\in\mathcal{L}_{\mathrm{H}}}\ell_s\le C_{\mathrm{H}},
  \qquad
  \lvert\mathcal{L}_{\mathrm{S}}\cup\mathcal{L}_{\mathrm{H}}\rvert
    \le N_{\mathrm{sess}}.
  \label{eq:cap}
\end{equation}
A policy maps $\mathcal{L}(t)$ to a victim on overflow and to
promotion and demotion candidates during a gap. The offline B\'el\'ady
oracle $v^{\star}$ evicts the session whose next request is farthest
in the future. An online policy $\pi$ may use only the causal
filtration $\mathcal{F}_{t}$, comprising the last gap, the turn index,
and the completion flag, and may use neither future arrivals nor agent-role
identity. The design target is to approximate $v^{\star}$ while
remaining a function of $\mathcal{F}_{t}$ alone. A learned policy
such as a small neural network could in principle fit the same
mapping, but the eviction decision sits on the critical path of every
cache miss and must complete within single-digit microseconds at the
hardware event rate, a latency regime that is incompatible with even a
minimal inference pass and that motivates the closed-form scoring
function developed below.

\subsection{\spear: Survival-Penalty Eviction for Agent Return-Gap}
\label{sec:esp}

The B\'el\'ady distance of a paused session depends on two quantities
that operate on different timescales, namely how long the current
tool gap will last and whether the session will return at all. \esp
approximates the first through a gap exponential moving average and
the second through a turn-indexed hazard, combining them into a
single eviction score that ranks all residents by estimated
next-reference distance.

\paragraph{Gap recurrence estimate}
The gap EMA tracks the characteristic return interval of session~$s$
with coefficient $\alpha{=}3/10$ as
\begin{equation}
  g_s
  \leftarrow
  \alpha\,g_{s,k}+(1-\alpha)\,g_s
  =\frac{3\,g_{s,k}+7\,g_s}{10}.
  \label{eq:ema}
\end{equation}
A session with a long smoothed gap is expected to remain idle longer,
so it should be evicted before a session that returns frequently.

\paragraph{Turn-indexed completion hazard}
A session closer to its final turn is more likely to complete and
release its KV permanently. The discrete hazard at turn $t$ is
\begin{equation}
  h(t)=\frac{d(t)}{n(t)},
  \qquad
  t\in\{0,\ldots,50\},
  \label{eq:hazard}
\end{equation}
where $n(t)$ sessions remain live at turn $t$ and $d(t)$ complete
there, counted over the observed population. The survival mass is the
local complement
\begin{equation}
  \sigma(t)
  =\max\bigl(\varepsilon,\,1-h(\min(t,50))\bigr),
  \qquad
  \varepsilon=0.01,
  \label{eq:surv}
\end{equation}
which requires a single table lookup, avoiding the sequential
dependence of the cumulative product-limit estimator.

\paragraph{Composite eviction score}
The two signals combine into
\begin{equation}
  \mathrm{score}(s)
  =
  \begin{cases}
    S_{\max},
      & s\text{ completed},\\[2pt]
    \dfrac{g_s}{\sigma(t)}+\bigl(1-\sigma(t)\bigr)P,
      & \text{otherwise},
  \end{cases}
  \label{eq:esp}
\end{equation}
where $P$ is a programmable completion penalty. A high score denotes
high eviction priority. The first term is the expected remaining wait
scaled by survival likelihood, so a session with low $\sigma$
near completion sees an inflated ratio that raises its eviction
priority. The
second term adds a penalty that grows as $\sigma$ falls, accelerating
eviction of sessions whose KV will soon have no future request. The
signal contribution and parameter sensitivity of~\eqref{eq:esp} are
evaluated in Sections~\ref{sec:sw-signal} and~\ref{sec:sw-main}.

Scoring determines who leaves the pool but not where a surviving
session should reside within the two-tier hierarchy. The
complementary placement problem is addressed next.

\subsection{\tide: Tiering in Idle-Window DMA Events}
\label{sec:okt}

A tool gap is wasted time from the perspective of the memory
hierarchy, because the paused session holds a tier slot but issues
no requests. \okt converts this idle window into a DMA budget and
migrates sessions between SRAM and HBM so that the lowest-scoring
residents occupy the fast tier when the next request arrives.

On a \texttt{gap\_start} event with estimated remaining duration
$\Delta$, the migration budget in tokens is
\begin{equation}
  \mathcal{B}=\Delta\cdot B,
  \label{eq:budget}
\end{equation}
where $B$ is the logical DMA bandwidth in tokens per nanosecond.
Let $\mathcal{H}$ be the HBM residents and $\mathcal{S}$ the SRAM
residents. Promotion and demotion then solve the complementary
selections
\begin{align}
  p
  &=\arg\min_{s\in\mathcal{H}}\mathrm{score}(s),
  \label{eq:promote}
  \\
  d
  &=\arg\max_{s\in\mathcal{S}}\mathrm{score}(s),
  \label{eq:demote}
\end{align}
subject to $\ell_p+\ell_d\le\mathcal{B}$ and to~\eqref{eq:cap} after
the swap. The key property of \okt is that it ranks migration
candidates with the same \esp score used for eviction, so eviction
quality directly determines placement quality, and the two mechanisms
share a single ranking signal rather than maintaining independent
state. This coupling is validated quantitatively in the component
ablation of Section~\ref{sec:sw-comp}.

\subsection{Joint Online Control Loop}
\label{sec:loop}

\esp and \okt share a single session register file updated on every
event. Algorithm~\ref{alg:unison} specifies the complete loop. Each
incoming event updates the register file, triggers an eviction scan
if capacity is violated, and initiates migrations if a gap is open
and the DMA engine is free. The two mechanisms read the same scores
from the same state, which is why a unified control plane
outperforms dual independent engines with periodic synchronization,
as the architectural ablation of Section~\ref{sec:arch}
demonstrates.

\begin{algorithm}[t]
\caption{\unison event loop}
\label{alg:unison}
\begin{algorithmic}[1]
\REQUIRE event $e$, live set $\mathcal{L}$, capacities
         $(C_{\mathrm{S}},C_{\mathrm{H}},N_{\mathrm{sess}})$,
         shared register file $R$, DMA state
\ENSURE  victim $v$ and logical DMA descriptors
\STATE write $e$ into $R$
         \COMMENT{turn, gap, completion}
\IF{$e$ is \texttt{gap\_end}}
  \STATE $g_s\leftarrow\alpha\,g_{s,k}+(1-\alpha)\,g_s$
         \COMMENT{\eqref{eq:ema}}
\ENDIF
\IF{$\mathcal{L}$ violates~\eqref{eq:cap}}
  \FOR{$s\in\mathcal{L}$}
    \STATE $\mathrm{score}(s)\leftarrow$~\eqref{eq:esp}
  \ENDFOR
  \STATE $v\leftarrow\arg\max_{s}\mathrm{score}(s)$
  \STATE evict $v$ from its tier
\ENDIF
\IF{$e$ is \texttt{gap\_start} and DMA is free}
  \STATE $\mathcal{B}\leftarrow\Delta\cdot B$
  \WHILE{$\mathcal{B}$ covers a legal pair}
    \STATE $p\leftarrow\arg\min_{s\in\mathcal{H}}\mathrm{score}(s)$
    \STATE $d\leftarrow\arg\max_{s\in\mathcal{S}}\mathrm{score}(s)$
    \STATE migrate $(p,d)$;
           $\mathcal{B}\leftarrow\mathcal{B}-\ell_p-\ell_d$
  \ENDWHILE
\ENDIF
\end{algorithmic}
\end{algorithm}

Section~\ref{sec:arch} maps this algorithmic contract to a
synthesizable fixed-point pipeline, and
Sections~\ref{sec:sweval}--\ref{sec:hweval} evaluate the resulting
policy against the B\'el\'ady oracle and competing baselines.
